% ----- CHAUPAI 26 -----

\newpage

\subsection*{\deva{षड्विंशतितम चौपाई} (Twenty-Sixth Chaupai)}
\addcontentsline{toc}{subsection}{\deva{षड्विंशतितम चौपाई} (Twenty-Sixth Chaupai) — \deva{संकट तें...}}

\begin{minipage}[t]{0.55\textwidth}
\vspace{0pt}

{\large \linespread{1.5}\selectfont
\deva{संकट तें हनुमान छुड़ावै।}\\
\deva{मन क्रम बचन ध्यान जो लावै॥}
\par}

\vspace{0.5em}

{\large \linespread{1.5}\selectfont
\telu{సంకట తేఁ హనుమాన ఛుడావై।}\\
\telu{మన క్రమ బచన ధ్యాన జో లావై॥}
\par}

\vspace{0.5em}

{\large \linespread{1.3}\selectfont
\textit{saṅkaṭa tẽ hanumāna chuṛāvai |}\\
\textit{mana krama bacana dhyāna jo lāvai ||}
\par}
\end{minipage}%
\hfill
\begin{minipage}[t]{0.42\textwidth}
\vspace{0pt}
\begin{tcolorbox}[anchorbox]
\textbf{The Perfect Alignment:} When your mind (\textit{Mana}), words (\textit{Bachan}), and actions (\textit{Krama}) are perfectly aligned toward a single goal, you can navigate your way out of any crisis (\textit{Sankat}). Hypocrisy and distraction are the true enemies of success.
\vspace{0.5em}\newline He turned a terrifying crisis around by perfectly aligning his sharp mind with his decisive, practical actions.\newline\null\hfill\href{https://www.valmiki.iitk.ac.in/sloka?field_kanda_tid=5&language=dv&field_sarga_value=53}{\textit{Sundara Kanda, 53:25-27}}
\end{tcolorbox}
\end{minipage}

\vspace{0.5em}
\subsubsection*{\textbf{Visual Rhythm Map:}}
\begin{table}[h!]
\centering
\renewcommand{\arraystretch}{1.2}
\setlength{\tabcolsep}{0pt}
\begin{tabularx}{\textwidth}{|*{16}{>{\centering\arraybackslash\hsize=1.0\hsize}X|}}
\hline
\multicolumn{4}{|>{\centering\arraybackslash\hsize=4.0\hsize}X|}{\textbf{\guru{सं}\laghu{क}\laghu{ट}} \par (4)} &
\multicolumn{2}{>{\centering\arraybackslash\hsize=2.0\hsize}X|}{\textbf{\guru{तें}} \par (2)} &
\multicolumn{5}{>{\centering\arraybackslash\hsize=5.0\hsize}X|}{\textbf{\laghu{ह}\laghu{नु}\guru{मा}\laghu{न}} \par (5)} &
\multicolumn{5}{>{\centering\arraybackslash\hsize=5.0\hsize}X|}{\textbf{\laghu{छु}\guru{ड़ा}\guru{वै}} \par (5)} \\
\hline
\end{tabularx}
\vspace{-1.5em}
\end{table}

\begin{table}[h!]
\centering
\renewcommand{\arraystretch}{1.2}
\setlength{\tabcolsep}{0pt}
\begin{tabularx}{\textwidth}{|*{16}{>{\centering\arraybackslash\hsize=1.0\hsize}X|}}
\hline
\multicolumn{2}{|>{\centering\arraybackslash\hsize=2.0\hsize}X|}{\textbf{\laghu{म}\laghu{न}} \par (2)} &
\multicolumn{2}{>{\centering\arraybackslash\hsize=2.0\hsize}X|}{\textbf{\laghu{क्र}\laghu{म}} \par (2)} &
\multicolumn{3}{>{\centering\arraybackslash\hsize=3.0\hsize}X|}{\textbf{\laghu{ब}\laghu{च}\laghu{न}} \par (3)} &
\multicolumn{3}{>{\centering\arraybackslash\hsize=3.0\hsize}X|}{\textbf{\guru{ध्या}\laghu{न}} \par (3)} &
\multicolumn{2}{>{\centering\arraybackslash\hsize=2.0\hsize}X|}{\textbf{\guru{जो}} \par (2)} &
\multicolumn{4}{>{\centering\arraybackslash\hsize=4.0\hsize}X|}{\textbf{\guru{ला}\guru{वै}} \par (4)} \\
\hline
\end{tabularx}
\vspace{-1.5em}
\end{table}

\vspace{1em}
\textbf{ \deva{सरल-संस्कृतम्}:}\\
 \deva{सङ्कटात् तमेव हनुमान् मोचयति}\\
 \deva{यः मनसा, कर्मणा, वचनेन च तस्य ध्यानं करोति॥}\\


Hanuman rescues that person from crisis and danger\\
who meditates on him with mind, action, and words/speech.


\vspace{1em}
\newpage
\textbf{\deva{प्रतिपदार्थः} (Meaning of each word):}
\begin{table}[h!]
\centering
\renewcommand{\arraystretch}{1.2}
\begin{tabularx}{\textwidth}{@{} l l X X @{}}
\toprule
\textbf{Awadhi} & \textbf{Sanskrit} & \textbf{English} & \textbf{Grammar \& Notes} \\
\midrule
\deva{संकट तें} / \telu{సంకట తేఁ} / \textit{saṅkaṭa tẽ}  &  \deva{सङ्कटात्}  &  From crisis  & Ablative postposition `-tẽ` \\
\deva{हनुमान} / \telu{హనుమాన} / \textit{hanumāna}  &  \deva{हनुमान्}  &  Hanuman  &   \\
\deva{छुड़ावै} / \telu{ఛుడావై} / \textit{chuṛāvai}  &  \deva{मोचयति}  &  Rescues/Frees  &   \\
\deva{मन} / \telu{మన} / \textit{mana}  &  \deva{मनसा}  &  With mind  &   \\
\deva{क्रम} / \telu{క్రమ} / \textit{krama}  &  \deva{कर्मणा}  &  With actions  & Metathesis of Karma \\
\deva{बचन} / \telu{బచన} / \textit{bacana}  &  \deva{वचनेन}  &  With words  & \\ 
\deva{ध्यान} / \telu{ధ్యాన} / \textit{dhyāna}  &  \deva{ध्यानम्}  &  Meditation  &   \\
\deva{जो लावै} / \telu{జో లావై} / \textit{jo lāvai}  &  \deva{यः करोति (आनयति)}  &  Whoever applies  &   \\
\bottomrule
\end{tabularx}
\end{table}



\newpage
